%=====================================================================
%  ANEXO B -- MARCO NORMATIVO: DESARROLLO COMPLETO
%  Responsable: Fabian Moreno Ugarte
%
%  GENERADO desde secciones/anexo_b_marco_normativo.md, que es la
%  FUENTE de este anexo. No editar aqui: editar el .md y reconvertir.
%
%  Desarrolla la Seccion 3, que en el cuerpo va comprimida. Las
%  remisiones internas (B.2, B.9...) van escritas en el .md: si se
%  renumera un subapartado, hay que corregirlas a mano.
%=====================================================================

\section{Marco normativo: desarrollo completo}
\label{anx:marco}

Este anexo recoge el desarrollo completo del marco normativo que la
Sección 3 presenta de forma resumida: el articulado citado literalmente,
el razonamiento que lleva de cada obligación a un requisito verificable
sobre el sistema y el contraste con el prototipo del equipo. Sus nueve
subapartados siguen el orden del análisis; la Sección 3 agrupa lo
esencial en cuatro. Como allí, las zonas en que la norma admite más de
una lectura razonable se presentan como riesgos a validar y no como
conclusiones, y la calificación jurídica definitiva corresponde al área
legal de Lima Airport Partners y, en su caso, a la Autoridad Nacional de
Protección de Datos Personales.

% ---- B.1 ----
\subsection{Marco normativo aplicable}

En el nivel constitucional, el artículo 2 de la Constitución de 1993
reconoce en su inciso 6 el derecho «a que los servicios informáticos,
computarizados o no, públicos o privados, no suministren informaciones
que afecten la intimidad personal y familiar», y en su inciso 7 el
derecho «al honor y a la buena reputación, a la intimidad personal y
familiar así como a la voz y a la imagen propias». De ambos se desprende
el criterio que atraviesa esta sección: captar la imagen de una persona
en un espacio abierto al público no es intrínsecamente ilícito, pero
constituye una injerencia en un derecho fundamental y debe superar un
examen de proporcionalidad.

La norma central es la Ley 29733, Ley de Protección de Datos Personales,
cuyo artículo 1 declara que su objeto es garantizar el derecho
fundamental previsto en el artículo 2 numeral 6 de la Constitución. Su
reglamento vigente fue aprobado por el Decreto Supremo 016-2024-JUS,
expedido el 29 de noviembre de 2024 y publicado al día siguiente, el 30
de noviembre de 2024, en la Separata Especial de \emph{El Peruano}. Su
Disposición Complementaria Derogatoria Única derogó el D.S.
003-2013-JUS. Su entrada en vigor se habría producido el 30 de marzo de
2025, dato que se enuncia con reserva: de los cuatro que se acaban de
citar es el único que no se ha contrastado todavía contra la fuente
oficial, y es precisamente el que fija el régimen transitorio, de modo
que debe confirmarse antes de apoyarse en él para fechar el paso de un
reglamento a otro. La consecuencia práctica es que los procedimientos,
cláusulas informativas y registros que LAP tenga construidos bajo el
reglamento anterior no son automáticamente conformes con el nuevo y
deben revisarse. \verificar{únicamente la fecha de entrada en vigor, el
30 de marzo de 2025, que es el dato que fija el régimen transitorio. La
expedición, la publicación y la disposición derogatoria sí se
contrastaron contra la separata de \emph{El Peruano} del 30 de noviembre
de 2024.}

El régimen se especifica para videovigilancia en la Directiva de
Tratamiento de Datos Personales mediante Sistemas de Videovigilancia,
identificada como Directiva 01-2020-JUS/DGTAIPD y aprobada por
Resolución Directoral 02-2020-JUS/DGTAIPD, publicada el 16 de enero de
2020. Sigue vigente y resulta aplicable, y esto ya no se afirma por
indicio sino por el texto de la norma: la Disposición Complementaria
Derogatoria Única del D.S. 016-2024-JUS, contrastada contra la separata
de \emph{El Peruano} del 30 de noviembre de 2024, deroga únicamente el
D.S. 003-2013-JUS, de modo que no hay derogación de la Directiva, ni
total ni parcial. La entrada en vigor del nuevo reglamento no la
desplazó, y otras entidades del Estado continúan invocándola como base
legal de sus propios sistemas con posterioridad a él.

A ese cuerpo se superpone un segundo régimen, de finalidad distinta, en
materia de seguridad ciudadana: el Decreto Legislativo 1218, que regula
el uso de las cámaras de videovigilancia, y la Ley 30120, Ley de Apoyo a
la Seguridad Ciudadana con Cámaras de Videovigilancia Públicas y
Privadas, ambos reglamentados por el Decreto Supremo 007-2020-IN, de 24
de abril de 2020. La denominación oficial completa y la fecha de
publicación del D.L. 1218 y de la Ley 30120 no se han contrastado
todavía contra su texto publicado, a diferencia del D.S. 007-2020-IN,
que sí se contrastó; ambas normas se citan, por tanto, por su
denominación usual, y esos dos datos deben confirmarse antes de
reproducirse en un documento dirigido al cliente. Este último importa
más de lo que el equipo había supuesto, porque su artículo 3 delimita el
ámbito de aplicación incluyendo a quienes son propietarios o
administradores de «establecimientos comerciales abiertos al público con
un aforo de cincuenta (50) personas o más». Un terminal aeroportuario
supera ese umbral con holgura, de modo que la aplicabilidad de este
régimen al Aeropuerto Internacional Jorge Chávez deja de ser hipótesis
remota y pasa a ser el supuesto que habría que descartar expresamente si
se pretendiera no aplicarlo. \verificar{denominación oficial completa y
fecha de publicación del D.L. 1218 y de la Ley 30120. El D.S.
007-2020-IN sí se contrastó contra el texto publicado en \emph{El
Peruano}.}

Queda así identificada la zona gris principal. Un terminal concesionado
a un operador privado no encaja limpiamente ni en la categoría de área
de dominio público ni en la de recinto privado cerrado: es un espacio de
titularidad y gestión privada, abierto al público, con funciones de
interés público y sujeto a regulación sectorial aeronáutica. De esa
calificación depende qué régimen resulta exigible y con qué intensidad.
Resolverla requeriría analizar el contrato de concesión y la normativa
aeronáutica, lo que excede el alcance del Capstone; es el primer punto a
elevar al área legal del cliente.

% ---- B.2 ----
\subsection{Naturaleza de los datos tratados}

Buena parte de la discusión de cumplimiento depende de una distinción
que conviene fijar con precisión, porque el sistema no trata un solo
tipo de dato sino cuatro, con regímenes distintos.

El primero es el fotograma. Contiene rostros identificables y constituye
sin discusión un dato personal en el sentido del artículo 2 numeral 4 de
la Ley 29733, que los define como «toda información sobre una persona
natural que la identifica o la hace identificable a través de medios que
pueden ser razonablemente utilizados». El criterio legal es el de
identificabilidad por medios razonables, no el de identificación
efectiva: que nadie mire la imagen no la convierte en dato no personal.

El segundo es el punto con coordenadas que producen los métodos basados
en puntos \cite{Song_2021_ICCV,Lin_2025_CVPR}: desligado del fotograma
del que se extrajo, no identifica por sí solo a una persona ni la hace
identificable por medios razonables.

El tercero es el dato espacio-temporal, que aparece en cuanto ese punto
se asocia a una marca de tiempo y a una zona del terminal. Aislado sigue
sin identificar a nadie; encadenado con otros del mismo individuo,
reconstruye una trayectoria, y una trayectoria dentro de un terminal
(dónde se dirigió alguien, cuánto permaneció, con quién coincidió) tiene
un poder identificante muy superior al de cualquiera de sus puntos por
separado. La diferencia entre el segundo tipo de dato y el tercero no es
algorítmica sino de política de persistencia: depende de si las salidas
se guardan y se asocian entre cuadros, no de qué red las produjo.

El cuarto es el dato biométrico. El artículo 2 numeral 5 define los
datos sensibles como los «datos personales constituidos por los datos
biométricos que por sí mismos pueden identificar al titular; datos
referidos al origen racial y étnico; ingresos económicos, opiniones o
convicciones políticas, religiosas, filosóficas o morales; afiliación
sindical; e información relacionada a la salud o a la vida sexual». La
fórmula es restrictiva y conviene leerla literalmente: no todo rasgo
corporal medido es dato sensible, sino el que por sí mismo permite
identificar al titular. Una plantilla facial cumple esa condición; un
mapa de densidad y un conjunto de coordenadas de cabezas, no.

De esta taxonomía se sigue un matiz que suele pasarse por alto: que la
salida del modelo no sea biométrica no significa que la entrada no lo
sea. El fotograma del que se derivan las coordenadas contiene rostros en
calidad suficiente para reconocimiento facial. Si se almacena, se
transmite fuera del perímetro o se conserva para reentrenamiento, el
análisis de riesgo cambia aunque el modelo desplegado nunca ejecute
reconocimiento facial.

Debe registrarse aquí una discrepancia terminológica detectada en el
material del equipo, que no es menor porque la ley trata ambos supuestos
como figuras jurídicas distintas. El material describe como «personas
anonimizadas» el resultado de un procesamiento que asigna a cada persona
un identificador persistente con el que se reconstruyen trayectorias. El
artículo 2 numeral 12 define la anonimización como el «tratamiento de
datos personales que impide la identificación o que no hace
identificable al titular de estos», y añade que «el procedimiento es
irreversible»; el numeral 13 define la disociación en términos casi
idénticos pero cierra con la nota opuesta: «el procedimiento es
reversible». Un identificador persistente que permite reagrupar las
apariciones de un mismo individuo no impide la identificación de forma
irreversible: es, en la terminología de la ley, una disociación, la
figura que el régimen europeo denomina seudonimización. La distinción
tiene consecuencias directas, porque el dato disociado sigue siendo dato
personal sujeto al régimen general, mientras que el anonimizado sale de
él.

La discrepancia no se limita al material de presentación. El módulo de
persistencia del prototipo incorporado al repositorio describe lo que
almacena como «solo posiciones, zonas, timestamps y el id temporal no
reversible», y su esquema define ese identificador como clave primaria
de una tabla de personas a la que se enlaza el historial completo de
posiciones. Un identificador así permite reagrupar las apariciones de
una misma persona y reconstruir su recorrido: el procedimiento es
reversible, y por tanto disociación en el sentido del numeral 13 y no
anonimización en el del numeral 12. Este informe no corrige la
terminología del equipo por su cuenta; deja registrada la discrepancia,
ahora en sus dos soportes (la presentación y el código), y la eleva como
punto abierto.

% ---- B.3 ----
\subsection{Base de legitimación y finalidad declarada}

El artículo 5 enuncia el principio de consentimiento, y el artículo 13.5
precisa que este «debe ser previo, informado, expreso e inequívoco». Ese
estándar es materialmente inalcanzable en un terminal: el pasajero no
elige ser observado, no puede negociar las condiciones y, en la
práctica, tampoco puede optar por un terminal sin cámaras. Un
consentimiento que no puede negarse sin renunciar al servicio no es
libre, de modo que apoyar el sistema en esa base sería construir la
legitimidad sobre una ficción.

La ley prevé la salida. Su artículo 14 enumera los supuestos en que no
se requiere consentimiento, y la Directiva 01-2020-JUS/DGTAIPD, en su
numeral 6.3, reconoce como bases del tratamiento por videovigilancia el
consentimiento, la autorización por ley o la concurrencia de alguno de
esos supuestos. El que la práctica invoca con más frecuencia es el
numeral 9, referido al tratamiento «necesario para salvaguardar
intereses legítimos del titular de datos personales por parte del
titular de datos personales o por el encargado de datos personales».
Conviene advertir que esta redacción no coincide con la del interés
legítimo del responsable propia del régimen europeo (el texto alude a
los intereses legítimos del titular de los datos), y que de esa
diferencia depende que el supuesto ampare o no una finalidad de gestión
operativa como la del proyecto. Qué base invoca LAP, y con qué
fundamento normativo, es información que solo el cliente puede aportar,
y sin ella el cartel del Anexo A no puede completarse.

Sobre la finalidad, en cambio, la norma no deja margen. El artículo 6
establece que los datos «deben ser recopilados para una finalidad
determinada, explícita y lícita» y que su tratamiento «no debe
extenderse a otra finalidad que no haya sido la establecida de manera
inequívoca como tal al momento de su recopilación», salvo actividades de
valor histórico, estadístico o científico cuando medie disociación o
anonimización. El artículo 28 numeral 4 lo repite como obligación: «no
utilizar los datos personales objeto de tratamiento para finalidades
distintas de aquellas que motivaron su recopilación, salvo que medie
procedimiento de anonimización o disociación».

De aquí se sigue una segunda discrepancia. El informe describe, entre
las aplicaciones posibles del seguimiento multicámara, la medición de la
permanencia de personas en zonas comerciales, y la vincula a un interés
comercial manifestado por el cliente. Medir permanencia en zona
comercial no es una variante de estimar aglomeraciones: es una finalidad
distinta. La primera busca saber cuántas personas hay en un punto en un
momento dado, para dimensionar la respuesta operativa; la segunda,
cuánto tiempo permanece una persona frente a una oferta comercial, lo
que constituye un análisis de comportamiento con valor económico propio.
Conforme a los artículos 6 y 28 numeral 4, requeriría su propia
declaración al momento de la recopilación, su propia base de
legitimación y su propio deber de información, y no podría ampararse en
la finalidad de seguridad ni en la de gestión de aglomeraciones. Este
informe deja la cuestión planteada sin resolverla, pero advierte que
presentar ambas finalidades como una sola sería precisamente el
deslizamiento que el artículo 6 prohíbe.

% ---- B.4 ----
\subsection{Principio de proporcionalidad y granularidad de salida}

El principio está consagrado en el artículo 7 de la Ley 29733, conforme
al cual «todo tratamiento de datos personales debe ser adecuado,
relevante y no excesivo a la finalidad para la que estos hubiesen sido
recopilados», y se reitera para videovigilancia en el numeral 6.4 de la
Directiva: «el tratamiento de los datos personales debe ser adecuado,
pertinente y no excesivo en relación con el ámbito y las finalidades
determinadas, legítimas y explícitas». Es el argumento más sólido a
favor del enfoque del proyecto y el eje de su defensa ante el cliente.

La estimación automática de densidad es idónea, porque detecta
aglomeraciones que la observación humana de decenas de monitores no
advierte a tiempo. En cuanto a la necesidad, que es el subprincipio
decisivo, si la finalidad declarada es estimar aglomeraciones, el medio
menos invasivo que la alcanza es un sistema que produce conteos
agregados y no identifica personas; uno con reconocimiento facial
lograría lo mismo siendo más invasivo, y por eso no superaría el examen.
La proporcionalidad en sentido estricto exige ponderar el beneficio
esperado (prevención de incidentes y mejora del flujo) contra la
injerencia en la imagen de miles de personas al día, y conviene señalar
aquí la debilidad del razonamiento en lugar de disimularla: esa
ponderación requiere evidencia sobre la magnitud real del problema en el
terminal (incidentes registrados, tiempos de cola, quejas) que solo LAP
puede aportar. Mientras no la aporte, el examen se sostiene sobre una
premisa no acreditada y así debe presentarse al cliente.

De este principio depende la decisión de granularidad de salida, la de
mayor consecuencia normativa del proyecto, y las tres opciones no son
equivalentes ante el artículo 7. El mapa de densidad de CSRNet
\cite{Li_2018_CVPR} produce una superficie agregada de la que el conteo
se obtiene por integración, sin posiciones individuales: es la opción
que menos datos trata para la finalidad declarada. La agregación por
zona añade una capa de reducción sobre cualquier método y produce el
mismo resultado operativo. La localización individual por puntos
\cite{Song_2021_ICCV,Lin_2025_CVPR} entrega la posición de cada persona,
que es más información de la que la finalidad requiere, y su licitud
depende entonces de que esas coordenadas no se persistan ni se
encadenen, es decir, de una garantía de política y no de arquitectura.
El criterio resultante es que, a igualdad de utilidad operativa, la
opción menos granular es la que supera el examen de necesidad, y
cualquier granularidad adicional debe justificarse por una utilidad que
la agregación no alcance. La elección del frente técnico no consta por
escrito al cerrar esta versión, y de ella depende que este apartado pase
de criterio a conclusión.

% ---- B.5 ----
\subsection{Privacidad desde el diseño y por defecto}

La privacidad desde el diseño no es en el régimen peruano una buena
práctica opcional, pero conviene enunciar con precisión de dónde se
desprende. El D.S. 016-2024-JUS no la nombra como principio autónomo: lo
que su Título Preliminar recoge, en el Artículo IX, son el principio de
transparencia (numeral 1) y el de responsabilidad proactiva (numeral 2),
conforme al cual el responsable no solo debe cumplir sino estar en
condiciones de acreditar que cumple, adoptando las medidas necesarias
antes de iniciar el tratamiento y no una vez iniciado. De ese segundo
principio se sigue la doble exigencia que este informe aplica y que la
literatura designa como privacidad por diseño (las garantías se
incorporan desde la concepción del tratamiento y no se añaden después) y
por defecto (la configuración más protectora debe venir activada de
fábrica, sin que nadie tenga que solicitarla). Importa señalar además
que el Artículo IX \textbf{no fija criterios ni contenidos mínimos de
acreditación}: enuncia el principio y deja al responsable la elección de
los medios con que lo demuestra. Las medidas que se proponen a
continuación no pueden, por tanto, contrastarse contra una lista
reglamentaria, porque no existe; se formulan como propuesta de este
frente y su suficiencia es, en último término, apreciación del área
legal de LAP.

La diferencia entre una y otra lectura no es retórica: si se tratara de
una recomendación, las medidas siguientes serían propuestas descartables
por conveniencia técnica; al ser exigencia normativa son requisitos, y
lo que queda abierto es cómo cumplir, no si cumplir. La dimensión por
defecto tiene además una consecuencia concreta: donde una opción admita
un ajuste más protector y otro menos protector (resolución de captura,
frecuencia de muestreo, granularidad de la salida, plazo de retención),
el valor de fábrica debe ser el más protector, y cualquier configuración
menos restrictiva debe ser una desviación deliberada, documentada y
justificada.

El principio solo es verificable si se traduce en medidas concretas. El
proyecto incorpora diez como requisitos del sistema: procesamiento en el
borde, de modo que el fotograma no salga del perímetro (M-01); no
persistencia del fotograma, que se procesa en memoria y se descarta
(M-02); salida agregada por zona, sin coordenadas individuales
persistidas (M-03); difuminado irreversible de rostros si algún
fotograma debe conservarse (M-04); plazo de retención definido con purga
automática (M-05); control de acceso por roles con registro de auditoría
(M-06); cifrado en tránsito y en reposo (M-07); carteles informativos
conforme al Anexo A (M-08); documentación del linaje de datos (M-09); y
supervisión humana, de modo que el sistema alerte y no ejecute acciones
automáticas sobre personas (M-10).

M-02 es la de mayor impacto sobre el perfil de riesgo, porque sin imagen
almacenada desaparece la mayor parte del riesgo de brecha, y también la
que más restringe al equipo técnico, ya que sin fotogramas conservados
no hay material propio para reentrenar ni para diagnosticar fallos. La
salida intermedia consiste en conservar un subconjunto reducido con
difuminado irreversible y plazo corto bajo procedimiento documentado.
Cuál de las dos versiones se propone al cliente queda abierto: la
estricta es más defendible jurídicamente, la atenuada es la única que
permite mejorar el modelo con datos del propio terminal, que es lo que
exige la brecha de dominio documentada en \cite{Lin_2025_CVPR}.

Estas diez medidas son requisitos, no una descripción del sistema
actual. Al cerrar esta versión, \textbf{M-03 no se cumple en el
prototipo incorporado al repositorio}, que persiste coordenadas
individuales asociadas a un identificador por persona; el contraste se
desarrolla en B.9. M-02, en cambio, sí se respeta. Conviene mantener la
distinción entre lo exigido y lo implementado, porque una lista de
medidas que se presente al cliente como estado alcanzado, cuando es
estado deseado, es exactamente el tipo de afirmación que este informe
evita.

% ---- B.6 ----
\subsection{Plazos de retención y supresión}

El fundamento está en el principio de calidad del artículo 8, que ordena
que los datos «deben conservarse de forma tal que se garantice su
seguridad y solo por el tiempo necesario para cumplir con la finalidad
del tratamiento». La obligación correlativa figura en el artículo 28
numeral 7: «suprimir los datos personales objeto de tratamiento cuando
hayan dejado de ser necesarios o pertinentes a la finalidad para la cual
hubiesen sido recopilados o hubiese vencido el plazo para su
tratamiento, salvo que medie procedimiento de anonimización o
disociación».

Para videovigilancia, la Directiva 01-2020-JUS/DGTAIPD concreta ese
criterio abierto en un plazo: sus numerales 6.13 y 7.18 fijan la
conservación de las imágenes en un rango de treinta a sesenta días como
máximo.

Aparece aquí una tensión normativa que el proyecto debe resolver por
diseño y que hasta ahora solo se había anticipado en abstracto. El
artículo 17.2 del D.S. 007-2020-IN obliga a los sujetos comprendidos en
su ámbito a «almacenar las imágenes, videos o audios grabados por un
plazo mínimo de cuarenta y cinco (45) días calendario, salvo disposición
distinta en normas sectoriales», y su artículo 17.1 impone entregar esa
información a la Policía Nacional del Perú o al Ministerio Público «en
un plazo máximo de veinticuatro (24) horas» cuando se presuma la
comisión de un delito o falta. El régimen de protección de datos empuja
hacia conservar lo mínimo y fija un techo de sesenta días; el de
seguridad ciudadana fija un piso de cuarenta y cinco. Ambos son
compatibles solo dentro de una ventana estrecha, y lo son únicamente
porque recaen sobre finalidades distintas.

La consecuencia de diseño es la que este frente ya venía recomendando y
que ahora queda respaldada con plazos concretos: tratar el sistema de
estimación de aglomeraciones como un subsistema separado del CCTV de
seguridad, con su propia finalidad, su propia base de legitimación y su
propio plazo de retención. Mezclarlos haría que el subsistema de conteo
heredara la obligación de conservación de cuarenta y cinco días propia
del video de seguridad, lo que perjudicaría el argumento de minimización
sin beneficio operativo alguno, dado que los conteos agregados no sirven
a la finalidad de seguridad ciudadana. Bajo M-02 y M-03, además, el
subsistema de conteo no almacena imágenes, de modo que la obligación del
artículo 17.2 recaería sobre el CCTV preexistente de LAP y no sobre lo
que el proyecto añade. El plazo concreto de retención de los datos
derivados no se fija aquí, porque depende de las necesidades analíticas
del cliente.

% ---- B.7 ----
\subsection{Derechos de los titulares e información al público}

El Título III de la Ley 29733 reconoce un catálogo de derechos que el
sistema debe permitir ejercer. El artículo 18 establece el derecho de
información, conforme al cual el titular debe ser informado «en forma
detallada, sencilla, expresa, inequívoca y de manera previa a su
recopilación» sobre la finalidad del tratamiento, los destinatarios, la
existencia del banco de datos y la identidad y domicilio de su titular,
las transferencias previstas, el tiempo durante el cual se conservarán
sus datos y la posibilidad de ejercer los derechos que la ley le
concede. El artículo 19 reconoce el derecho de acceso; el 20, los de
actualización, inclusión, rectificación y supresión; el 21, el de
impedir el suministro; el 22, el de oposición; el 24, el de tutela ante
la Autoridad Nacional o mediante hábeas data; y el 25, el de ser
indemnizado.

Merece mención aparte el artículo 23, que consagra el derecho al
tratamiento objetivo: el derecho del titular «a no verse sometido a una
decisión con efectos jurídicos sobre él o que le afecte de manera
significativa, sustentada únicamente en un tratamiento de datos
personales destinado a evaluar determinados aspectos de su personalidad
o conducta». Es el fundamento normativo directo de M-10: mientras el
sistema se limite a emitir alertas agregadas que un operador humano
interpreta, no produce decisiones automatizadas sobre personas
individuales y el artículo no se activa; si sus salidas alimentaran una
medida automática sobre un pasajero concreto, sí lo haría.

En videovigilancia, el deber de información se cumple de forma
característica mediante señalética. El numeral 6.11 de la Directiva fija
el contenido mínimo del cartel (identidad y domicilio del responsable,
forma de ejercer los derechos y lugar donde obtener la información del
artículo 18) y establece una dimensión mínima de 297 × 210 milímetros;
el Anexo A propone un modelo construido sobre esos elementos. Dado que
el consentimiento no es la base habilitante, la legitimidad del sistema
descansa enteramente en que supere el examen de proporcionalidad y en
que este deber se cumpla de manera efectiva: el cartel es, junto con la
proporcionalidad, uno de los dos pilares de la licitud del tratamiento,
y no un trámite.

Varios de sus campos no pueden completarse con información que el equipo
posea: la razón social registral y el domicilio del responsable, la base
de legitimación invocada, el canal de atención de derechos y el plazo de
conservación son datos que solo LAP puede proporcionar. Un cartel con
esos campos en blanco no cumple el deber de información, de modo que su
obtención condiciona la medida M-08.

% ---- B.8 ----
\subsection{Encargados de tratamiento y flujos hacia terceros}

El artículo 2 numeral 6 define al encargado del banco de datos
personales como quien «realiza el tratamiento de los datos personales
por encargo del titular del banco de datos personales». Cualquier
proveedor que trate datos por cuenta de LAP (etiquetado, plataformas de
MLOps, almacenamiento de objetos, infraestructura de inferencia) ocupa
esa posición y queda sujeto al artículo 30, conforme al cual los datos
tratados por cuenta de terceros «no pueden aplicarse o utilizarse con un
fin distinto al que figura en el contrato o convenio celebrado ni ser
transferidos a otras personas, ni aun para su conservación», y deben
suprimirse una vez ejecutada la prestación salvo autorización expresa.

Cuando el destinatario se sitúa fuera del país entra en juego el flujo
transfronterizo, definido en el artículo 2 numeral 8 y regulado en el
artículo 15, que solo lo permite «si el país destinatario mantiene
niveles de protección adecuados conforme a la presente Ley» y, en su
defecto, obliga al emisor a garantizar que el tratamiento se efectúe
conforme a la ley peruana. El artículo 11 recoge el principio
correlativo de nivel de protección adecuado. El D.S. 016-2024-JUS
contempla además un alcance extraterritorial del régimen peruano, fijado
en el Artículo VI de su Título Preliminar. De sus criterios de conexión,
dos son los que este proyecto debe examinar: el numeral 3.1, que atrae
al régimen peruano el tratamiento vinculado a la oferta de bienes o
servicios a titulares situados en territorio peruano, y el numeral 3.2,
que hace lo propio con el análisis del comportamiento de esos titulares
y la elaboración de perfiles. Ninguno de los dos puede descartarse de
antemano si alguna etapa del procesamiento llegara a ejecutarse fuera
del país; cuál de ellos resulta aplicable a un encargado contratado por
LAP es calificación que corresponde a su área legal.

De aquí se derivan tres consecuencias sobre decisiones que el frente
técnico probablemente ya está tomando. Si el entrenamiento o la
inferencia se ejecutan en nube alojada fuera del Perú, existe un flujo
transfronterizo que debe analizarse. El uso de un modelo preentrenado
descargado del exterior no constituye por sí mismo transferencia de
datos personales, porque no se exportan datos; enviar fotogramas del
terminal a una interfaz externa para inferencia sí lo constituiría. Y el
servidor externo de evaluación de NWPU-Crowd, al que se remiten
predicciones porque el conjunto no libera las etiquetas de prueba,
implica un envío a un tercero fuera del Perú: en fase académica el
riesgo es acotado, porque se envían predicciones sobre imágenes de un
conjunto público, pero remitir por esa vía cualquier artefacto derivado
de material del terminal constituiría una transferencia internacional
con todas sus consecuencias.

La decisión entre inferencia \emph{on-premise} e inferencia en nube
externa es, a juicio de este frente, la de mayor impacto legal del
proyecto, porque determina si existe o no transferencia internacional.
Al cerrar esta versión no consta por escrito, y la medida M-01 la
anticipa en el sentido más protector sin que ello equivalga a que esté
tomada.

% ---- B.9 ----
\subsection{Auditoría: ¿el sistema trata datos biométricos?}

Esta subsección contrasta las decisiones técnicas del equipo contra los
requisitos anteriores. La auditoría se formula de manera condicional
sobre los enfoques revisados en la Sección 2 porque la elección de
arquitectura no consta todavía en la documentación del proyecto, pero
conviene precisar que \textbf{sí consta en código}: el frente técnico
incorporó al repositorio un prototipo de detección, seguimiento y fusión
entre cámaras, y lo que ese prototipo implementa se examina al final de
este apartado. La pregunta que decide el régimen aplicable no es si hay
cámaras, sino si en algún punto del procesamiento se extrae de la imagen
una característica que por sí sola permita identificar a una persona, en
el sentido del artículo 2 numeral 5.

Bajo CSRNet \cite{Li_2018_CVPR} ninguna etapa construye una plantilla
biométrica y la salida, un mapa de densidad, es agregada por
construcción, de modo que no admite reidentificación posterior. Respecto
de este enfoque es razonable afirmar que el tratamiento no recae sobre
datos sensibles, con la salvedad de que el fotograma de entrada sí
contiene rostros.

Bajo P2PNet o P2R \cite{Song_2021_ICCV,Lin_2025_CVPR} tampoco se
construyen plantillas faciales, de modo que el tratamiento no sería
biométrico, pero la conformidad queda condicionada a la política de
persistencia: si las coordenadas se almacenan y se encadenan entre
cuadros, el sistema se aproxima a un seguimiento de trayectorias, que
como se expuso en B.3 es una finalidad distinta de la declarada. La
conclusión es favorable, condicionada a que M-03 se implemente
efectivamente.

Bajo el enfoque de seguimiento multicámara la calificación cambia de
naturaleza, y conviene decirlo sin atenuación. Los trabajos de esta
familia revisados en la Sección 2 asignan un identificador global a cada
persona y lo mantienen entre cámaras mediante características de
apariencia y módulos de reidentificación. Un vector de apariencia
construido para volver a reconocer a un individuo en otra cámara es, por
su función, una característica extraída del cuerpo de la persona cuyo
propósito es individualizarla. Si permitiera por sí mismo identificar al
titular, quedaría comprendido en el artículo 2 numeral 5 y el
tratamiento recaería sobre datos sensibles, con el régimen reforzado que
ello implica; si no lo permitiera, seguiría siendo dato personal sujeto
al régimen general, pero el sistema estaría reconstruyendo trayectorias
individuales, es decir, tratando el tercer tipo de dato descrito en B.2
en su forma de mayor poder identificante. En ninguna de las dos lecturas
el resultado es neutro, y en ambas excede lo que la finalidad de estimar
aglomeraciones requiere conforme al artículo 7. La calificación concreta
no puede cerrarse sin conocer qué se computa y qué se persiste, y
corresponde al área legal de LAP.

Conviene ahora contrastar ese análisis con lo que el prototipo del
equipo implementa realmente, porque el resultado corrige en un punto
importante lo que este frente había anticipado y agrava en otro.

\textbf{Sobre la reidentificación por apariencia, la preocupación no se
confirma.} El material de presentación del equipo contemplaba
reidentificación mediante \emph{embeddings} de apariencia, lo que habría
chocado frontalmente con el requisito no funcional RNF-01, que prohíbe
derivar rasgos que permitan la identificación indirecta de las personas.
El texto literal de ese requisito, y el documento en que consta, no
figuran en el repositorio al cerrar esta versión, de modo que su
formulación se recoge aquí según lo reportado por este frente y no según
su fuente, y la conformidad con él que se afirma a continuación queda
sujeta a ese contraste. El prototipo no hace eso. La fusión entre
cámaras se resuelve como asociación espacio-temporal sobre un plano
común compartido (dos cámaras que observan la misma explanada desde
ángulos distintos), y esa capa es la obligatoria. Solo ante ambigüedad
entre varios candidatos próximos interviene un descriptor de desempate
compuesto por un histograma de color del torso y la proporción
alto/ancho de la caja, que el propio código documenta como no biométrico
y que, según su diseño, nunca se emplea como identificador único; ese
descriptor tampoco se persiste, pues vive únicamente en memoria mientras
la persona está activa. Un histograma de color de la ropa no es un rasgo
del cuerpo y cambia cuando la persona se cambia de abrigo: no constituye
dato biométrico en el sentido del artículo 2 numeral 5, y es
sustancialmente menos identificante que un \emph{embedding} de
reidentificación. La decisión de diseño va, por tanto, en la dirección
que el RNF-01 exige, y así debe reconocerse. Queda por confirmar que el
criterio se sostenga si más adelante se busca mayor exactitud en la
fusión, porque es precisamente ahí donde reaparece la tentación del
\emph{embedding}. \verificar{texto literal del RNF-01 y documento en que
consta. No figura en el repositorio al cerrar esta versión, de modo que
su formulación se recoge según lo reportado por este frente y no según
su fuente.}

\textbf{Sobre lo que se persiste, en cambio, aparece un hallazgo no
conforme que no estaba previsto.} El esquema de base de datos del
prototipo almacena una tabla de personas, con un identificador global
como clave primaria y las marcas de primera y última detección, y una
tabla de posiciones que guarda, por cada detección, el identificador de
la persona, la cámara, las coordenadas, el instante y la zona. Eso es
exactamente el tercer tipo de dato descrito en B.2, es decir, el dato
espacio-temporal encadenado a un identificador persistente, y está
escrito en disco, no procesado en memoria. La medida M-03, que exige
salida agregada por zona sin coordenadas individuales persistidas, no se
cumple en la implementación actual. La medida M-02 sí se respeta, porque
no se almacenan fotogramas. La consecuencia jurídica es la desarrollada
en B.4: una trayectoria individual persistida excede lo que la finalidad
de estimar aglomeraciones requiere conforme al artículo 7, y su licitud
ya no puede apoyarse en que la salida sea agregada, porque no lo es.

Conviene señalar además que el comentario del propio módulo de
persistencia describe lo almacenado como «el id temporal no reversible».
Esa calificación reproduce en el código la discrepancia terminológica
registrada en B.2: un identificador que es clave primaria de un
historial completo de posiciones permite reagrupar y reconstruir el
recorrido de una persona, de modo que el procedimiento es reversible en
el sentido del artículo 2 numeral 13 y constituye disociación, no
anonimización. La corrección no es cosmética, porque de ella depende si
el dato sale o no del régimen general de la ley: no sale.

Del contraste resultan cuatro hallazgos no conformes de naturaleza
distinta. El primero es la ausencia de la Evaluación de Impacto en
Protección de Datos, exigible antes de iniciar tratamientos de alto
riesgo, supuesto en el que la videovigilancia masiva figura
expresamente; es el más grave en términos formales, por tratarse de una
obligación incumplida y no de una carencia de evidencia, pero también el
más fácil de subsanar, ya que el material sustantivo está redactado en
la Sección 4 y lo que falta es formalizarlo. El segundo es el
incumplimiento de M-03 por persistencia de coordenadas individuales,
descrito en el párrafo anterior; a diferencia de los demás, \textbf{es
un defecto de implementación y no una ausencia de evidencia}, lo que lo
hace subsanable por el propio equipo y sin depender de terceros, ya sea
agregando por zona antes de escribir o fijando una retención muy corta
sobre la tabla de posiciones. El tercero es la ausencia de validación de
exactitud en el dominio de destino, pues no existe validación en el
terminal y la brecha de dominio documentada en \cite{Lin_2025_CVPR}
impide extrapolar las cifras publicadas. El cuarto es la ausencia de
evaluación de sesgo demográfico, que ninguno de los trabajos revisados
reporta de forma desagregada. Los dos últimos no son defectos de
implementación sino ausencias de evidencia, y ambos requieren acceso a
material del terminal. Quedan además pendientes, por depender de
decisiones abiertas, la política de retención, la localización del
procesamiento y la trazabilidad de licencias de los conjuntos de datos.
